\documentclass[11pt]{article}
\usepackage[margin=1in]{geometry}
\usepackage{amsmath,amssymb,amsthm}
\usepackage{graphicx}
\usepackage[hidelinks]{hyperref}

\newtheorem{theorem}{Theorem}
\newtheorem{lemma}{Lemma}

\newcommand{\Lf}{L^1([0,1],X)}
\newcommand{\Linf}{L^\infty([0,1],X^*)}
\newcommand{\Sf}{S_f}
\newcommand{\weak}{\mathrm{weak}}

\title{Weak Compactness of State Spaces in \(L^1([0,1],X)\)}
\author{run fa\_banach\_001, agent\_lane\_16}
\date{2026-06-28}

\begin{document}
\maketitle

\paragraph{Status.}
\texttt{full\_solution\_likely\_valid} for Problem 3.18 of arXiv:2509.12780,
under the same separable-dual hypotheses used in the source paper's
Theorems 2.3 and 2.5. Problem 3.19, the nonseparable-dual variant, is not
claimed.

\section{Source Problem}

The source is Soumitra Daptari and Saurabh Dwivedi, ``A study on state spaces
in classical Banach spaces,'' arXiv:2509.12780, to appear in Colloquium
Mathematicum \cite{DaptariDwivedi2025}. For \(x\in S(X)\), the state space is
\[
  S_x=\{x^*\in S(X^*): x^*(x)=1\}.
\]
Under the standing identification \(L^1([0,1],X)^*=L^\infty([0,1],X^*)\) when
\(X^*\) is separable, the source defines \(S_f\) analogously for
\(f\in S(L^1([0,1],X))\).

Problem 3.18, on page 18, asks for the weak compactness analogue of Theorem
2.5:

\begin{center}
\includegraphics[width=0.96\textwidth]{figures/open_problem_crop.png}
\end{center}

In local source coordinates this is
\texttt{data/parsed/arxiv\_sources/2509.12780/Banach\_CM\_.tex}, lines
442--447.

\section{Idea of the Proof}

The norm-compactness theorem in the source proves that a nonsmooth fibre set
of positive measure lets one choose many separated pointwise states. For weak
compactness, pairwise norm separation is not enough: weakly compact sets can
contain separated sequences. The extra ingredient is to alternate between two
measurably selected pointwise states by Rademacher signs. After evaluation at
one fixed vector of \(X\), this produces the usual Rademacher sequence in
\(L^\infty\). The closed span of that sequence is isometric to \(\ell_1\), so
no subsequence is weakly Cauchy. Eberlein--Smulian then forbids weak
compactness of \(S_f\).

\section{Theorem}

\begin{theorem}
Let \(X\) be a real or complex Banach space such that \(X^*\) is separable,
and let \(\mu\) be Lebesgue measure on \([0,1]\). For
\(f\in S(L^1(\mu,X))\), regard
\[
  S_f=\{h\in S(L^\infty(\mu,X^*)):\int_0^1 h(t)(f(t))\,dt=1\}.
\]
Then \(S_f\) is weakly compact in \(L^\infty(\mu,X^*)\) if and only if \(f\)
is a smooth point of \(L^1(\mu,X)\).

Equivalently, by Theorem 2.3 of \cite{DaptariDwivedi2025}, \(S_f\) is weakly
compact if and only if \(f(t)\ne0\) almost everywhere and
\[
  \frac{f(t)}{\|f(t)\|}
\]
is a smooth point of \(X\) for almost every \(t\).
\end{theorem}

\begin{proof}
If \(f\) is smooth, then \(S_f\) is a singleton, hence weakly compact. We prove
the converse.

We shall use the following elementary observation repeatedly. If \(h\in S_f\),
then the equality
\[
  1=\int_0^1 h(t)(f(t))\,dt\le \int_0^1 \|h(t)\|\,\|f(t)\|\,dt
     \le \int_0^1\|f(t)\|\,dt=1
\]
forces \(h(t)(f(t))=\|f(t)\|\) for almost every \(t\), after taking real
parts in the complex case. Thus, off the zero set of \(f\), \(h(t)\) is a
pointwise state of \(f(t)/\|f(t)\|\) almost everywhere.

First suppose that \(Z=\{t:\|f(t)\|=0\}\) has positive measure. Choose
\(h_0\in S_f\), choose \(y\in S(X)\), and choose \(y^*\in S(X^*)\) with
\(y^*(y)=1\). Since \(Z\) is a positive-measure non-atomic Lebesgue set, it
supports Rademacher functions \(r_n:Z\to\{-1,1\}\) such that every finite sign
pattern occurs on a subset of positive measure. Define
\[
 h_n(t)=
 \begin{cases}
   r_n(t)y^*, & t\in Z,\\
   h_0(t), & t\notin Z.
 \end{cases}
\]
Then \(h_n\in S_f\) for every \(n\). The bounded evaluation map
\[
  P:L^\infty(\mu,X^*)\longrightarrow L^\infty(Z),\qquad
  P(H)(t)=H(t)(y),
\]
sends \(h_n\) to \(r_n\). If \(S_f\) were weakly compact, Eberlein--Smulian
would give a weakly convergent subsequence of \((h_n)\), and hence a weakly
convergent subsequence of \((r_n)\) in \(L^\infty(Z)\). This is impossible:
for every subsequence \((r_{n_j})\), the closed span of \(\{r_{n_j}\}\) is
isometric to \(\ell_1\), since for finite scalars \((a_j)\),
\[
 \left\|\sum_j a_j r_{n_j}\right\|_\infty=\sum_j |a_j|.
\]
The unit vector basis of \(\ell_1\) has no weakly Cauchy subsequence: after
passing to any subsequence, the bounded functional with alternating
coordinates \(1,-1,1,-1,\ldots\) prevents convergence. This contradiction
shows that \(\mu(Z)=0\).

Now assume \(f(t)\ne0\) almost everywhere and suppose that \(f\) is not smooth
in \(L^1(\mu,X)\). By Theorem 2.3 of \cite{DaptariDwivedi2025}, the set
\[
  E=\left\{t:\frac{f(t)}{\|f(t)\|}\text{ is not smooth in }X\right\}
\]
has positive measure. Put \(u_t=f(t)/\|f(t)\|\) on \(E\). Since
\(X^*\) is separable, \(X\) is separable and \(S(X^*)\), with the norm
topology, is Polish. For some \(\delta>0\), the set of \(t\in E\) for which
\(\operatorname{diam} S_{u_t}\ge \delta\) has positive measure. The measurable
graph argument used in the source paper, together with von Neumann selection,
gives measurable maps \(\varphi,\psi:E_0\to S(X^*)\) on a positive-measure
subset \(E_0\subseteq E\) such that
\[
  \varphi(t),\psi(t)\in S_{u_t},\qquad
  \|\varphi(t)-\psi(t)\|\ge\delta
\]
for all \(t\in E_0\).

Let \((x_m)\) be a countable dense subset of \(S(X)\). After replacing
\(E_0\) by a positive-measure subset \(B\), we may fix \(x_0=x_m\) such that
\[
  a(t):=(\varphi(t)-\psi(t))(x_0)
\]
satisfies \(|a(t)|\ge \delta/2\) on \(B\). Let \(r_n:B\to\{-1,1\}\) be
Rademacher functions on \(B\), and let \(h_0\in S_f\). Define
\[
 h_n(t)=
 \begin{cases}
   \frac{1+r_n(t)}2\,\varphi(t)+\frac{1-r_n(t)}2\,\psi(t), & t\in B,\\
   h_0(t), & t\notin B.
 \end{cases}
\]
On \(B\), each value is either \(\varphi(t)\) or \(\psi(t)\), hence is a
state of \(u_t\). Outside \(B\), \(h_0\) is already a state for \(f\) almost
everywhere. Thus \(h_n\in S_f\) for all \(n\).

Define a bounded linear map
\[
  Q:L^\infty(\mu,X^*)\longrightarrow L^\infty(B),\qquad
  Q(H)(t)=\frac{2H(t)(x_0)}{a(t)}.
\]
It is bounded because \(|a(t)|\ge\delta/2\). If
\(m(t)=(\varphi(t)+\psi(t))/2\) on \(B\), then
\[
  Q(h_n)(t)=Q(m)(t)+r_n(t)\qquad(t\in B).
\]
Weak compactness of \(S_f\) would again imply, by Eberlein--Smulian and
boundedness of \(Q\), that some subsequence of \((r_n)\) is weakly convergent
in \(L^\infty(B)\), contradicting the same \(\ell_1\)-basis argument. Hence
\(f\) must be smooth.
\end{proof}

\section{Consequence for Problem 3.18}

Combining the theorem above with Theorems 2.3 and 2.5 of
\cite{DaptariDwivedi2025} gives the exact weak compactness analogue requested
in Problem 3.18:
\[
  S_f\text{ weakly compact}
  \quad\Longleftrightarrow\quad
  S_f\text{ norm compact}
  \quad\Longleftrightarrow\quad
  f\text{ is smooth}.
\]
Equivalently, under the source hypotheses, this happens precisely when
\(f(t)\ne0\) a.e. and \(f(t)/\|f(t)\|\) is smooth in \(X\) a.e.

\section{Verification and Novelty Check}

The proof is purely functional-analytic; no computational verification is
needed. The key points for human review are:
\begin{itemize}
\item the measurable selection of two uniformly separated fibre states on a
positive-measure subset;
\item the bounded evaluation map \(Q\), whose denominator is bounded away from
zero after choosing a fixed vector \(x_0\);
\item the use of Eberlein--Smulian and the isometric \(\ell_1\) span of
Rademacher functions in \(L^\infty\).
\end{itemize}

Bounded novelty checks were performed in the run indexes for
\texttt{2509.12780}, state spaces, weak compactness, \(L^1([0,1],X)\), and
\(S_f\). No prior packet for this problem was found. Web searches for the
exact paper title and phrases such as ``state spaces weak compactness
\(L^1\)'' and ``\(S_f\) weakly compact \(L^1([0,1],X)\)'' found the source
arXiv record but no later answer.

\section{Limitations}

This packet solves Problem 3.18 in the separable-dual framework of the source
paper. It does not address Problem 3.19, where \(X^*\) is not separable and
the measurable-selection machinery used both here and in the source paper is
not directly available.

\begin{thebibliography}{9}
\bibitem{DaptariDwivedi2025}
S. Daptari and S. Dwivedi,
``A study on state spaces in classical Banach spaces,''
arXiv:2509.12780, 2025; to appear in Colloquium Mathematicum.

\bibitem{DiestelUhl1977}
J. Diestel and J. J. Uhl,
\emph{Vector Measures},
Mathematical Surveys, No. 15, American Mathematical Society, 1977.

\bibitem{Srivastava1998}
S. M. Srivastava,
\emph{A Course on Borel Sets},
Graduate Texts in Mathematics 180, Springer, 1998.
\end{thebibliography}

\end{document}
